\section{Data}\label{sec:data}

\paragraph{Treatment.}
STR prohibition petitions and executions come from Chicago Data Portal disclosures with coordinates; \texttt{TractProhibitionDatesProcessor} maps each building footprint to Census tracts, producing tract-level onset months that feed \texttt{TreatmentIndicatorProcessor}.
The analytic sample retains Cook County tract boundaries clipped to municipal borders when mapping.

\paragraph{Outcome.}
Monthly rent levels originate from Zip-code Zillow Observed Rent observations (\texttt{Zip\_zori\_uc\_sfrcondomfr\_sm\_month.csv}) expanded to tract-months via the ZIP--tract crosswalk from \texttt{ZipTractCrosswalkProcessor} plus \texttt{TimeSeriesZipToTractProcessor}.
Because conversion relies on overlaps between ZIP polygons and tract polygons, heterogeneous housing supply within oversized ZIP footprints injects smoothing relative to hypothetical unit-level rents.

\paragraph{Covariates.}
American Community Survey tract covariates (median income/house value, renters' share, age, attainment, baseline rent summaries) arrive through \texttt{CensusDataLoader} and merge through \texttt{DIDCovariateProcessor}, aligning with ACS vintage documented in-code.

\paragraph{Matched sample.}
\texttt{TrendMatchingProcessor} pairs treated tracts to nearest-sloped control neighbors in pre-ban rent dynamics (elastic net distance with $k$-nearest neighbors defaults catalogued in \texttt{docs/DID\_CS\_RUNBOOK.md}).
Descriptive exports (\texttt{did\_descriptive\_*.csv}) and panel excerpts (\texttt{did\_panel\_data.csv}) document coverage for slides or appendices.
